% PhD Thesis / Master's Thesis Template
% Compatible with most universities (adjust as needed)
% Use with the book class by default (frontmatter/mainmatter are supported).

\documentclass[12pt,oneside]{book}
\usepackage[utf8]{inputenc}
\usepackage[T1]{fontenc}
\usepackage{mathptmx}  % Times New Roman
\usepackage{graphicx}
\usepackage{appendix}
\usepackage{listings}
\usepackage{color}
\usepackage{amsmath,amssymb,amsfonts}
\usepackage{hyperref}
\usepackage{algorithm}
\usepackage{algorithmicx}
\usepackage{algpseudocode}
\usepackage{amsthm}
\usepackage{setspace}
\usepackage{geometry}
\usepackage{titlesec}
\usepackage{afterpage}
\usepackage{natbib}

% Page margins (adjust for your university)
\geometry{
  letterpaper,
  left=1.5in,
  right=1in,
  top=1in,
  bottom=1in
}

% Hyperlink setup
\hypersetup{
  colorlinks=true,
  linkcolor=blue,
  citecolor=blue,
  urlcolor=blue
}

% Theorem environments (optional)
\theoremstyle{plain}
\newtheorem{theorem}{Theorem}[chapter]
\newtheorem{lemma}[theorem]{Lemma}
\newtheorem{proposition}[theorem]{Proposition}
\newtheorem{corollary}[theorem]{Corollary}
\theoremstyle{definition}
\newtheorem{definition}{Definition}[chapter]

% Chapter title format
\titleformat{\chapter}
  {\normalfont\fontsize{14pt}{\baselineskip}\bfseries}
  {\thechapter}{20pt}{\fontsize{14pt}{\baselineskip}\bfseries}

\titleformat{\section}
  {\normalfont\fontsize{12pt}{\baselineskip}\bfseries}
  {\thesection}{15pt}{\fontsize{12pt}{\baselineskip}\bfseries}

\titleformat{\subsection}
  {\normalfont\fontsize{12pt}{\baselineskip}\itshape}
  {\thesubsection}{15pt}{\fontsize{12pt}{\baselineskip}\itshape}

% Spacing
\doublespacing
%\onehalfspacing  % Use if preferred

% Title Page
\title{
  \fontsize{16pt}{\baselineskip}\textbf{Your Thesis Title: A Comprehensive Study}\\[2cm]
  \fontsize{14pt}{\baselineskip}A Dissertation Submitted to the Graduate Faculty\\
  in Partial Fulfillment of the Requirements for the Degree of\\[0.5cm]
  \fontsize{14pt}{\baselineskip}\textbf{Doctor of Philosophy}\\[0.5cm]
  \fontsize{12pt}{\baselineskip}in\\[0.5cm]
  \fontsize{12pt}{\baselineskip}\textbf{Computer Science}
}

\author{
  \fontsize{14pt}{\baselineskip}\textbf{By}\\[0.5cm]
  \fontsize{14pt}{\baselineskip}\textbf{Your Name}\\[2cm]
}

\date{
  \fontsize{12pt}{\baselineskip}\textbf{Year: 2024}\\[1cm]
  \fontsize{12pt}{\baselineskip}Department of Computer Science\\
  University Name\\
  City, State, Country
}

\begin{document}

% Front matter
\frontmatter
\pagenumbering{roman}

% Title
\maketitle

% Copyright
\thispagestyle{empty}
\vspace*{\fill}
\begin{center}
Copyright \copyright\ Your Name, 2024\\
All Rights Reserved
\end{center}
\vspace*{\fill}
\clearpage

% Committee
\thispagestyle{empty}
\vspace*{3cm}
\begin{center}
\textbf{APPROVAL OF THE DISSERTATION COMMITTEE}
\end{center}
\vspace{1cm}

\begin{center}
This dissertation has been accepted by the Graduate Faculty of the University in partial fulfillment of the requirements for the degree of Doctor of Philosophy.
\end{center}

\vspace{2cm}

\begin{center}
\begin{tabular}{lr}
\hline
Dr. Committee Chair Name & \hspace{2in} \\
Chair of Committee & Date \\
\end{tabular}
\end{center}

\vspace{1cm}

\begin{center}
\begin{tabular}{lr}
\hline
Dr. Committee Member Name & \hspace{2in} \\
Committee Member & Date \\
\end{tabular}
\end{center}

\clearpage

% Abstract
\begin{center}
\textbf{ABSTRACT}\\[1cm]
\textbf{Your Name}\\
\textbf{Your Thesis Title: A Comprehensive Study}\\[0.5cm]
\textbf{Doctor of Philosophy in Computer Science}\\[0.5cm]
\textbf{University Name, 2024}
\end{center}

\vspace{1cm}

\begin{singlespace}
Your abstract goes here. It should be 150-350 words summarizing:
\begin{itemize}
  \item The problem being addressed
  \item The methodology used
  \item The main findings/results
  \item The contributions and implications
\end{itemize}
Keep it concise and informative. This is what people read to decide if they want to read the full thesis.
\end{singlespace}

\clearpage

% Acknowledgments
\chapter*{Acknowledgments}

I would like to express my gratitude to:

\begin{itemize}
  \item My advisor, Dr. Advisor Name, for guidance and support
  \item My committee members for valuable feedback
  \item My colleagues and collaborators
  \item My family and friends for their encouragement
  \item Funding agencies (if applicable)
\end{itemize}

\clearpage

% Table of Contents
\tableofcontents

% List of Tables
\listoftables

% List of Figures
\listoffigures

% List of Algorithms
\listofalgorithms

% List of Abbreviations
\chapter*{List of Abbreviations}
\begin{tabular}{ll}
AI & Artificial Intelligence \\
ML & Machine Learning \\
DL & Deep Learning \\
CNN & Convolutional Neural Network \\
RNN & Recurrent Neural Network \\
NLP & Natural Language Processing \\
API & Application Programming Interface \\
\end{tabular}

\clearpage

% Main body
\mainmatter
\pagenumbering{arabic}

\chapter{Introduction}
\label{chapter:introduction}

\section{Background and Motivation}

Provide context for your research. Why is this problem important? 
What are the real-world applications?

Example text: Recent advances in machine learning have led to 
significant breakthroughs in various domains \cite{smith2023}. 
However, challenges remain in [specific area] \cite{jones2022}.

\section{Problem Statement}

Clearly define the problem your research addresses. Be specific 
and precise about what you are solving.

\section{Research Objectives}

List your main research objectives:

\begin{enumerate}
  \item Objective 1
  \item Objective 2
  \item Objective 3
\end{enumerate}

\section{Contributions}

Summarize the key contributions of this dissertation:

\begin{enumerate}
  \item A novel framework for [contribution 1]
  \item Theoretical analysis proving [contribution 2]
  \item Extensive empirical validation demonstrating [contribution 3]
\end{enumerate}

\section{Organization}

This dissertation is organized as follows. Chapter~\ref{chapter:literature} 
reviews related work. Chapter~\ref{chapter:methodology} describes the 
proposed methodology. Chapter~\ref{chapter:results} presents experimental 
results. Chapter~\ref{chapter:discussion} discusses implications and 
limitations. Chapter~\ref{chapter:conclusion} concludes with a summary 
and future directions.

\chapter{Literature Review}
\label{chapter:literature}

This chapter reviews relevant literature and establishes the 
theoretical foundation for this research.

\section{Theoretical Background}

Provide necessary background theory and concepts.

\section{Related Work}

Review existing approaches and methods. Organize thematically:

\subsection{Approach A}

[Review of Approach A]

\subsection{Approach B}

[Review of Approach B]

\section{Research Gap}

Identify what is missing in existing work and how this dissertation 
addresses that gap.

\chapter{Methodology}
\label{chapter:methodology}

This chapter describes your research methodology.

\section{System Overview}

Provide an overview of your approach. Include figures showing 
the overall architecture.

\section{Component Description}

Describe each component in detail.

\section{Theoretical Analysis}

Present any theoretical guarantees or analysis.

\section{Algorithm}

Present your algorithm (see Algorithm~\ref{alg:main}).

\begin{algorithm}
\caption{Algorithm Name}
\label{alg:main}
\begin{algorithmic}
\State \textbf{Input:} [Input description]
\State \textbf{Output:} [Output description]
\State Initialize parameters
\For{each iteration}
  \State [Step description]
\EndFor
\State \textbf{return} [Final result]
\end{algorithmic}
\end{algorithm}

\chapter{Results}
\label{chapter:results}

This chapter presents experimental results.

\section{Experimental Setup}

Describe datasets, baselines, metrics, and implementation details.

\section{Main Results}

Present main results in tables and figures.

\begin{table}[htbp]
\caption{Performance comparison.}
\label{tab:results}
\begin{center}
\begin{tabular}{|l|c|c|c|}
\hline
Method & Metric 1 & Metric 2 & Metric 3 \\
\hline
Baseline 1 & 85.2 & 84.8 & 85.5 \\
Baseline 2 & 87.1 & 86.5 & 87.8 \\
Our Approach & \textbf{91.3} & \textbf{90.9} & \textbf{91.5} \\
\hline
\end{tabular}
\end{center}
\end{table}

\section{Ablation Studies}

Analyze contribution of each component.

\section{Statistical Analysis}

Include statistical significance tests where appropriate.

\chapter{Discussion}
\label{chapter:discussion}

\section{Interpretation of Results}

Explain what the results mean.

\section{Implications}

Discuss broader implications of findings.

\section{Limitations}

Acknowledge limitations of your work:

\begin{itemize}
  \item Limitation 1
  \item Limitation 2
  \item Limitation 3
\end{itemize}

\section{Future Work}

Suggest directions for future research.

\chapter{Conclusion}
\label{chapter:conclusion}

Summarize the dissertation:

\begin{itemize}
  \item What problem was addressed?
  \item What was the approach?
  \item What were the main findings?
  \item What are the contributions?
  \item What are the future directions?
\end{itemize}

% Bibliography
\bibliographystyle{plainnat}
\bibliography{references}

% Appendices
\appendix

\chapter{Additional Material}

Include supplementary material here.

\section{Extra Tables}

\section{Extra Figures}

\section{Code}

Include pseudocode or key implementation details.

\end{document}
